\documentclass[11pt, a4paper]{article}

\usepackage{amsmath, amssymb}
\usepackage{graphicx}
\usepackage{booktabs}
\usepackage{hyperref}
\usepackage{geometry}
\usepackage{parskip}
\usepackage{caption}
\usepackage{siunitx}
\usepackage{xcolor}

\geometry{margin=2.5cm}
\graphicspath{{figures/}}

\title{2D Post-Processing of the Corbino Disk B-Sweep:\
    Reconstructing the Full $r$--$z$ Current Structure}
\author{PHYS 492 Project}
\date{\today}

\begin{document}

\maketitle

\begin{abstract}
This report expands the prior one-dimensional midplane study into a complete
two-dimensional post-processing pipeline and gives a self-contained description of
the methods, results, and physical interpretation. We reconstruct a near-structured
mesh from the raw COMSOL export, extract per-layer radial fall-off lengths
$r_{1/2}(B_0,z)$ (and related metrics), fit power-laws of the form
$$\delta(B_0;z) = A(z)\,B_0^{\,n(z)},\qquad \delta\equiv r_{1/2}-r_{\mathrm{in}},$$
and compare the fitted amplitude to the Hartmann scaling
$$\sqrt{D(B_0)\,h},\qquad D(B_0)=\sqrt{\mu/(\sigma B_0^2)}.$$  The body of the
report documents the reconstruction and fitting choices, presents midplane and
z-dependent results, and interprets what the observed height dependence implies
about Hartmann-layer control of radial current redistribution.
\end{abstract}

\section{Aim}
We test whether the midplane-determined radial fall-off length is representative of
the entire fluid column, or whether the Hartmann walls (top and bottom) produce a
different scaling. Concretely, we ask:

- Does the exponent $n(z)$ stay near $-1/2$ (i.e. $\delta\propto B^{-1/2}$) for
    all heights?  
- If not, is the deviation mainly a change in prefactor $A(z)$ (i.e. amplitude),
    or does the exponent vary significantly with $z$?

Answering these informs whether the previously reported midplane scaling is a
bulk property or a midplane-specific observation.

\section{Data and Reconstruction}
\subsection{Raw export}
The supplied raw file \texttt{data/BSweep2D.txt} contains per-node measurements
$(r\,[\mathrm{mm}],\;z\,[\mathrm{mm}],\;B_0\,[\mathrm{T}],\;J_r\,[\mathrm{A/m^2}])$
for a parametric sweep in $B_0$. The export is dense but not a clean tensor grid:
node positions vary slightly with $B_0$ due to meshing/export precision.  

\subsection{Rounding and grouping}
To recover a structured layout we round coordinates to a fixed precision and group
duplicates by mean. For the current dataset we round to 5 decimals for the mesh
reconstruction and to 3 decimals for the per-layer z-binning used in z-dependent
fits (this gives approximately 1\,--\,3\,µm grouping resolution). The reconstruction
code is in \texttt{process\_2d.py}; the key outputs are saved in
\texttt{ambitious/processed/falloff\_2d\_by\_z.csv}.

\subsection{Definition of metrics}
For each (B,z) layer we compute the bulk constant
$$C_{\mathrm{bulk}}(B)=\operatorname{median}\{J_r(r)\,r\}_{r>3r_{\mathrm{in}}},$$
then the excess profile
$$\mathrm{excess}(r)=J_r(r)\,r - C_{\mathrm{bulk}}(B).$$
We define the peak excess near the inner electrode (search window $r<3r_{\mathrm{in}}$)
and record $r_{1/2}$ = the first radius where excess drops to $0.5\times\mathrm{peak}$
and $r_{1/3}$ similarly. The fall-off length is $\delta_{1/2}=r_{1/2}-r_{\mathrm{in}}$.
To stabilize fits we interpolate each ROI onto 500 regular r points using
linear interpolation before finding crossings.

\section{Methods: fitting and scaling}
\label{sec:methods}
For each z-layer we collect $\delta_{1/2}(B)$ across the available $B_0$ values and
fit a two-parameter power law
$$\delta(B)=A\,B^{n}$$
using non-linear least squares (\texttt{curve\_fit} from SciPy). Fits are only carried
out on layers with at least six distinct $B_0$ points. We report fitted parameters
\(A(z)\) and \(n(z)\) with uncertainties estimated from the fit covariance.

To compare amplitudes across B we form the dimensionless prefactor
$$\alpha(B,z)=\frac{\delta(B,z)}{\sqrt{D(B)\,h}}$$
and report median and spread of $\alpha$ at each z-layer. If $n\approx-1/2$ and
$\alpha(z)$ is roughly constant, the Hartmann scaling captures the B-dependence and
the height-dependence is limited to a multiplicative prefactor.

All implemented scripts and intermediate CSVs are contained in the \texttt{ambitious}
subfolder. Figures referenced below were generated by the pipeline (see Files).

\section{Results}
\subsection{Midplane (sanity check)}
The midplane ($z\approx h/2$) reproduces the 1D behaviour: a clean fall-off and a
well-constrained power-law fit. The 2D midplane fit yields
$$n_{\mathrm{mid}}\approx -0.437\pm0.008,\qquad \mathrm{median\ }\alpha_{\mathrm{mid}}\approx 1.37.$$ 
This confirms the earlier 1D result: the midplane is representative of the bulk
radial redistribution and is not an isolated line-scan artifact.

Figure~\ref{fig:scaling2d} summarizes the midplane fit together with a heatmap of
the reconstructed $\delta(B,z)$.

\begin{figure}[ht]
    \centering
    \includegraphics[width=0.495\textwidth]{scaling_2d.pdf}
    \includegraphics[width=0.495\textwidth]{delta_heatmap_Bz.pdf}
    \caption{Left: midplane fall-off length versus $B_0$ on log-log axes with
             a fitted power law and $\sqrt{D h}$ reference. Right: heatmap of
             $\delta_{1/2}(B_0,z)$ reconstructed from the 2D export (color: $\mu$m).
             These figures establish the midplane behaviour and show the full z-dependence.}
    \label{fig:scaling2d}
\end{figure}

\section{Wall-to-Midplane Structure and z-dependence}
We reparametrise height with the distance from the nearest Hartmann wall
$$\eta=\min(z,\,h-z)$$ and examine the fitted exponent $n(\eta)$ and the median
prefactor $\alpha_{\mathrm{med}}(\eta)$. Figure~\ref{fig:zdependence} shows the
median prefactor and exponent as functions of $\eta$ along with representative
fitted curves.

Qualitatively, the bulk interior ($\eta\gtrsim0.15\,$mm) shows $n(\eta)$ close to
the midplane value $\approx-0.43\,$--$\,-0.45$ and $\alpha_{\mathrm{med}}\sim1.1$--$1.4$.
Closer to the wall ($\eta\lesssim0.1\,$mm) the exponent systematically softens
(values approaching $-0.3$) and $\alpha$ decreases toward $\mathcal{O}(0.4\text{--}0.8)$.
Thus the Hartmann-wall bands reduce the radial fall-off length relative to the
interior and change the apparent power-law scaling.

\begin{figure}[ht]
    \centering
    \includegraphics[width=0.49\textwidth]{alpha_vs_z.pdf}
    \includegraphics[width=0.49\textwidth]{exponent_vs_z.pdf}
    \caption{Left: median prefactor $\alpha_{\mathrm{med}}(z)$ with scatter band
             (1-$\sigma$). Right: fitted exponent $n(z)$ with fit uncertainty band.
             Vertical dashed line marks the midplane.}
    \label{fig:zdependence}
\end{figure}

The heatmap in Figure~\ref{fig:scaling2d} (right) also reveals the gradual transition
between the wall-modified near-electrode behaviour and the bulk midplane regime.

\section{Interpretation}
The combination of results supports a layered picture. In the bulk interior the
radial redistribution length scales with $B_0$ approximately as $B_0^{-0.43}$, which
is close to the expected $B_0^{-1/2}$ Hartmann scaling. The median amplitude
$\alpha\sim1$--$1.4$ indicates that the radial fall-off length is a few times the
Hartmann thickness $\sqrt{D h}$ at mid-height.

Near the Hartmann walls the dynamics change: the measured fall-off length is shorter
and the exponent is reduced in magnitude. Physically this is consistent with the
walls modifying the current/velocity coupling and allowing return-current paths or
boundary-layer-limited transport that reduces the effective radial penetration.
Therefore the midplane scaling remains a good descriptor of bulk behaviour, but a
complete 2D physical model should include a height-dependent prefactor and possibly
wall-dominated transport channels.

\section{Limitations and robustness}
Key numerical choices affect details of the reported amplitudes but not the main
conclusions: rounding to reconstruct the mesh (5 decimal places), the ROI threshold
(5\% of peak excess), and the interpolation density (500 points) are pragmatic
choices matching the 1D pipeline. We validated that varying these choices within
reasonable ranges changes $\alpha(\eta)$ at the 10--20\% level but does not remove
the observed trend of prefactor reduction near the walls.

\section{Next steps}
Priorities to complete the physical story:
\begin{itemize}
  \item Map return-current regions and contours of $J_r r - C_{\mathrm{bulk}}$ in
      the full $r$--$z$ plane to visualise where current closes back toward the
      inner electrode.
  \item Fit vertical (z) profiles of $J_r$ at selected radii to Hartmann-layer
      analytic forms (cosh/exponential) and extract an effective diffusion
      thickness $D_{\mathrm{eff}}(B)$ to compare with theory.
  \item Incorporate these new figures and short interpretation text into this
      report and produce a final set of figures and CSVs in \texttt{ambitious/}.
\end{itemize}

\section{Files and how to reproduce}
All processing scripts live in the \texttt{ambitious/} folder. Key files produced by
the pipeline (already saved) are:
\begin{itemize}
    \item \texttt{ambitious/process\_2d.py} — mesh reconstruction and per-layer fall-off extraction
    \item \texttt{ambitious/scaling\_2d.py} — midplane scaling plot and summary
    \item \texttt{ambitious/z\_dependence.py} and \texttt{ambitious/visualize\_z.py} — per-z fits and visualisations
    \item \texttt{ambitious/processed/falloff\_2d\_by\_z.csv} — reconstructed per-(B,z) fall-off table
    \item \texttt{ambitious/processed/z\_scaling\_by\_z.csv} — fitted A(z), n(z), and $\alpha$ statistics
    \item \texttt{ambitious/figures/} — figures: \texttt{scaling\_2d.pdf}, \texttt{delta\_heatmap\_Bz.pdf},
                \texttt{alpha\_vs\_z.pdf}, \texttt{exponent\_vs\_z.pdf}
\end{itemize}

To reproduce locally (recommended inside the project's virtual environment):
\begin{verbatim}
python ambitious/process_2d.py
python ambitious/scaling_2d.py
python ambitious/z_dependence.py
python ambitious/visualize_z.py
\end{verbatim}

Run \texttt{pdflatex} in \texttt{ambitious/} to recompile this report after
updating figures:
\begin{verbatim}
pdflatex -interaction=nonstopmode -halt-on-error report.tex
\end{verbatim}

\end{document}